% Studies-in-Algebra-and-Group-Theory - Lecture 10: Modules and the Classification of Finitely Generated Abelian Groups
% Author: S. D. V. Stephens
% Date: 2026-03-10
% Compile with: pdflatex -shell-escape

\documentclass{report}
\input{~/university/preamble.tex}
% preamble-darkmode.tex
% Dark mode extension for lecture notes
% Usage: % preamble-darkmode.tex
% Dark mode extension for lecture notes
% Usage: % preamble-darkmode.tex
% Dark mode extension for lecture notes
% Usage: \input{~/university/preamble-darkmode.tex} AFTER preamble.tex
%
% This provides:
% - Dark background with light text
% - Material Design color palette
% - Ocean blue gradient colors (p1-p9)
% - Enhanced TikZ/PGFPlots for dark backgrounds
% - qq tcolorbox environment
% - tkz-euclide for geometric constructions

% ===========================================
% DARK MODE PACKAGE
% ===========================================
\usepackage{darkmode}
\enabledarkmode

% ===========================================
% ADDITIONAL PACKAGES
% ===========================================
\usepackage{changepage}
\usepackage{ulem}
\usepackage{colortbl}
\usepackage{tkz-euclide}
\usepackage{capt-of}

% ===========================================
% EXTENDED TIKZ LIBRARIES
% ===========================================
\usetikzlibrary{patterns,3d,backgrounds}
\usepgfplotslibrary{groupplots}

% Upgrade pgfplots compatibility
\pgfplotsset{compat=1.18}

% ===========================================
% OCEAN BLUE GRADIENT (p1-p9)
% ===========================================
\definecolor{p1}{HTML}{caf0f8}
\definecolor{p2}{HTML}{ade8f4}
\definecolor{p3}{HTML}{90e0ef}
\definecolor{p4}{HTML}{48cae4}
\definecolor{p5}{HTML}{00b4d8}
\definecolor{p6}{HTML}{0096c7}
\definecolor{p7}{HTML}{0077b6}
\definecolor{p8}{HTML}{023e8a}
\definecolor{p9}{HTML}{03045e}

% Dark background color
\definecolor{pag}{HTML}{293133}

% ===========================================
% MATERIAL DESIGN COLORS
% ===========================================
% Note: myred, myyellow already defined in main preamble
% These are additional/alternate versions
\definecolor{matred}{HTML}{f44336}
\definecolor{matpink}{HTML}{e81e63}
\definecolor{matpurple}{HTML}{9c27b0}
\definecolor{matdeeppurple}{HTML}{673ab7}
\definecolor{matindigo}{HTML}{3f51b5}
\definecolor{matblue}{HTML}{2196f3}
\definecolor{matlightblue}{HTML}{03a9f4}
\definecolor{matcyan}{HTML}{00bcd4}
\definecolor{matteal}{HTML}{009688}
\definecolor{matgreen}{HTML}{4caf50}
\definecolor{matlightgreen}{HTML}{8bc34a}
\definecolor{matlime}{HTML}{cddc39}
\definecolor{matyellow}{HTML}{ffeb3b}
\definecolor{matamber}{HTML}{ffc107}
\definecolor{matorange}{HTML}{ff9800}
\definecolor{matdeeporange}{HTML}{ff5722}
\definecolor{matbrown}{HTML}{795548}
\definecolor{matgray}{HTML}{9e9e9e}
\definecolor{matbluegray}{HTML}{607d8b}

% ===========================================
% COLORLET FOR TIKZ
% ===========================================
\colorlet{xcol}{blue!60!black}

% ===========================================
% DARK MODE TCOLORBOX
% ===========================================
\newtcolorbox{qq}[2][]{%
    boxrule=0.75pt,
    sharp corners,
    colframe=white,
    colback=pag,
    coltext=white,
    #1,
}

% Dark definition box
\newtcolorbox{darkdfn}[2][]{%
    enhanced,
    boxrule=1pt,
    sharp corners,
    colframe=p5,
    colback=pag,
    coltext=white,
    coltitle=white,
    fonttitle=\bfseries,
    title=#2,
    #1,
}

% Dark theorem box
\newtcolorbox{darkthm}[2][]{%
    enhanced,
    boxrule=1pt,
    sharp corners,
    colframe=matpurple,
    colback=pag,
    coltext=white,
    coltitle=white,
    fonttitle=\bfseries,
    title=#2,
    #1,
}

% ===========================================
% TIKZ 3D FIX
% ===========================================
% Small fix for canvas is xy plane at z
% https://tex.stackexchange.com/a/48776/121799
\makeatletter
\tikzoption{canvas is xy plane at z}[]{%
    \def\tikz@plane@origin{\pgfpointxyz{0}{0}{#1}}%
    \def\tikz@plane@x{\pgfpointxyz{1}{0}{#1}}%
    \def\tikz@plane@y{\pgfpointxyz{0}{1}{#1}}%
    \tikz@canvas@is@plane}
\makeatother

% ===========================================
% CONDITIONAL LINE TIKZSET
% ===========================================
\tikzset{
    conditional line/.style 2 args={
        /utils/exec={\pgfmathsetmacro{\xcoordA}{#1}},
        /utils/exec={\pgfmathsetmacro{\xcoordB}{#2}},
        insert path={
            \ifdim\xcoordA pt>\xcoordB pt
            (\xcoordA,0) -- (\xcoordB,0)
            \fi
        },
    },
}

% ===========================================
% PGFPLOTS DARK MODE DEFAULTS
% ===========================================
\pgfplotsset{
    dark axis/.style={
        axis line style={white},
        tick style={white},
        label style={white},
        grid style={pag!70},
        every axis label/.append style={white},
        every tick label/.append style={white},
    },
}

% ===========================================
% TIKZ EXTERNALIZATION (for fast compilation)
% ===========================================
% This creates .md5 cache files for figures
% Requires: pdflatex -shell-escape
%
% To enable, add AFTER loading this file:
%   \enabletikzexternalize
%
\newcommand{\enabletikzexternalize}{%
  \usetikzlibrary{external}%
  \tikzexternalize[prefix=figures/]%
}

% ===========================================
% CONVENIENT SHORTCUTS FOR DARK PLOTS
% ===========================================
\newcommand{\darkplot}[1]{%
    \begin{tikzpicture}
    \begin{axis}[
        dark axis,
        grid=both,
        major grid style={line width=.2pt,draw=pag!70},
        #1
    ]
}


% ===========================================
% QUICK GEOMETRY MACROS (tkz-euclide)
% ===========================================
% Draw a point: \point{name}{x}{y}
\newcommand{\gpoint}[3]{\tkzDefPoint(#2,#3){#1}}

% Draw line through points: \gline{A}{B}
\newcommand{\gline}[2]{\tkzDrawLine[color=p4](#1,#2)}

% Draw circle: \gcircle{center}{radius}
\newcommand{\gcircle}[2]{\tkzDrawCircle[color=p5](#1,#2)}

% Mark angle: \gangle{A}{B}{C}
\newcommand{\gangle}[3]{\tkzMarkAngle[color=p3,size=0.5](#1,#2,#3)}

% ===========================================
% USAGE EXAMPLE
% ===========================================
% In your document:
%
% \begin{tikzpicture}
%   \begin{axis}[dark axis, ...]
%     \addplot[p4, thick] {sin(deg(x))};
%   \end{axis}
% \end{tikzpicture}
%
% Or with qq box:
% \begin{qq}{My Title}
%   Content here in dark box
% \end{qq}
 AFTER preamble.tex
%
% This provides:
% - Dark background with light text
% - Material Design color palette
% - Ocean blue gradient colors (p1-p9)
% - Enhanced TikZ/PGFPlots for dark backgrounds
% - qq tcolorbox environment
% - tkz-euclide for geometric constructions

% ===========================================
% DARK MODE PACKAGE
% ===========================================
\usepackage{darkmode}
\enabledarkmode

% ===========================================
% ADDITIONAL PACKAGES
% ===========================================
\usepackage{changepage}
\usepackage{ulem}
\usepackage{colortbl}
\usepackage{tkz-euclide}
\usepackage{capt-of}

% ===========================================
% EXTENDED TIKZ LIBRARIES
% ===========================================
\usetikzlibrary{patterns,3d,backgrounds}
\usepgfplotslibrary{groupplots}

% Upgrade pgfplots compatibility
\pgfplotsset{compat=1.18}

% ===========================================
% OCEAN BLUE GRADIENT (p1-p9)
% ===========================================
\definecolor{p1}{HTML}{caf0f8}
\definecolor{p2}{HTML}{ade8f4}
\definecolor{p3}{HTML}{90e0ef}
\definecolor{p4}{HTML}{48cae4}
\definecolor{p5}{HTML}{00b4d8}
\definecolor{p6}{HTML}{0096c7}
\definecolor{p7}{HTML}{0077b6}
\definecolor{p8}{HTML}{023e8a}
\definecolor{p9}{HTML}{03045e}

% Dark background color
\definecolor{pag}{HTML}{293133}

% ===========================================
% MATERIAL DESIGN COLORS
% ===========================================
% Note: myred, myyellow already defined in main preamble
% These are additional/alternate versions
\definecolor{matred}{HTML}{f44336}
\definecolor{matpink}{HTML}{e81e63}
\definecolor{matpurple}{HTML}{9c27b0}
\definecolor{matdeeppurple}{HTML}{673ab7}
\definecolor{matindigo}{HTML}{3f51b5}
\definecolor{matblue}{HTML}{2196f3}
\definecolor{matlightblue}{HTML}{03a9f4}
\definecolor{matcyan}{HTML}{00bcd4}
\definecolor{matteal}{HTML}{009688}
\definecolor{matgreen}{HTML}{4caf50}
\definecolor{matlightgreen}{HTML}{8bc34a}
\definecolor{matlime}{HTML}{cddc39}
\definecolor{matyellow}{HTML}{ffeb3b}
\definecolor{matamber}{HTML}{ffc107}
\definecolor{matorange}{HTML}{ff9800}
\definecolor{matdeeporange}{HTML}{ff5722}
\definecolor{matbrown}{HTML}{795548}
\definecolor{matgray}{HTML}{9e9e9e}
\definecolor{matbluegray}{HTML}{607d8b}

% ===========================================
% COLORLET FOR TIKZ
% ===========================================
\colorlet{xcol}{blue!60!black}

% ===========================================
% DARK MODE TCOLORBOX
% ===========================================
\newtcolorbox{qq}[2][]{%
    boxrule=0.75pt,
    sharp corners,
    colframe=white,
    colback=pag,
    coltext=white,
    #1,
}

% Dark definition box
\newtcolorbox{darkdfn}[2][]{%
    enhanced,
    boxrule=1pt,
    sharp corners,
    colframe=p5,
    colback=pag,
    coltext=white,
    coltitle=white,
    fonttitle=\bfseries,
    title=#2,
    #1,
}

% Dark theorem box
\newtcolorbox{darkthm}[2][]{%
    enhanced,
    boxrule=1pt,
    sharp corners,
    colframe=matpurple,
    colback=pag,
    coltext=white,
    coltitle=white,
    fonttitle=\bfseries,
    title=#2,
    #1,
}

% ===========================================
% TIKZ 3D FIX
% ===========================================
% Small fix for canvas is xy plane at z
% https://tex.stackexchange.com/a/48776/121799
\makeatletter
\tikzoption{canvas is xy plane at z}[]{%
    \def\tikz@plane@origin{\pgfpointxyz{0}{0}{#1}}%
    \def\tikz@plane@x{\pgfpointxyz{1}{0}{#1}}%
    \def\tikz@plane@y{\pgfpointxyz{0}{1}{#1}}%
    \tikz@canvas@is@plane}
\makeatother

% ===========================================
% CONDITIONAL LINE TIKZSET
% ===========================================
\tikzset{
    conditional line/.style 2 args={
        /utils/exec={\pgfmathsetmacro{\xcoordA}{#1}},
        /utils/exec={\pgfmathsetmacro{\xcoordB}{#2}},
        insert path={
            \ifdim\xcoordA pt>\xcoordB pt
            (\xcoordA,0) -- (\xcoordB,0)
            \fi
        },
    },
}

% ===========================================
% PGFPLOTS DARK MODE DEFAULTS
% ===========================================
\pgfplotsset{
    dark axis/.style={
        axis line style={white},
        tick style={white},
        label style={white},
        grid style={pag!70},
        every axis label/.append style={white},
        every tick label/.append style={white},
    },
}

% ===========================================
% TIKZ EXTERNALIZATION (for fast compilation)
% ===========================================
% This creates .md5 cache files for figures
% Requires: pdflatex -shell-escape
%
% To enable, add AFTER loading this file:
%   \enabletikzexternalize
%
\newcommand{\enabletikzexternalize}{%
  \usetikzlibrary{external}%
  \tikzexternalize[prefix=figures/]%
}

% ===========================================
% CONVENIENT SHORTCUTS FOR DARK PLOTS
% ===========================================
\newcommand{\darkplot}[1]{%
    \begin{tikzpicture}
    \begin{axis}[
        dark axis,
        grid=both,
        major grid style={line width=.2pt,draw=pag!70},
        #1
    ]
}


% ===========================================
% QUICK GEOMETRY MACROS (tkz-euclide)
% ===========================================
% Draw a point: \point{name}{x}{y}
\newcommand{\gpoint}[3]{\tkzDefPoint(#2,#3){#1}}

% Draw line through points: \gline{A}{B}
\newcommand{\gline}[2]{\tkzDrawLine[color=p4](#1,#2)}

% Draw circle: \gcircle{center}{radius}
\newcommand{\gcircle}[2]{\tkzDrawCircle[color=p5](#1,#2)}

% Mark angle: \gangle{A}{B}{C}
\newcommand{\gangle}[3]{\tkzMarkAngle[color=p3,size=0.5](#1,#2,#3)}

% ===========================================
% USAGE EXAMPLE
% ===========================================
% In your document:
%
% \begin{tikzpicture}
%   \begin{axis}[dark axis, ...]
%     \addplot[p4, thick] {sin(deg(x))};
%   \end{axis}
% \end{tikzpicture}
%
% Or with qq box:
% \begin{qq}{My Title}
%   Content here in dark box
% \end{qq}
 AFTER preamble.tex
%
% This provides:
% - Dark background with light text
% - Material Design color palette
% - Ocean blue gradient colors (p1-p9)
% - Enhanced TikZ/PGFPlots for dark backgrounds
% - qq tcolorbox environment
% - tkz-euclide for geometric constructions

% ===========================================
% DARK MODE PACKAGE
% ===========================================
\usepackage{darkmode}
\enabledarkmode

% ===========================================
% ADDITIONAL PACKAGES
% ===========================================
\usepackage{changepage}
\usepackage{ulem}
\usepackage{colortbl}
\usepackage{tkz-euclide}
\usepackage{capt-of}

% ===========================================
% EXTENDED TIKZ LIBRARIES
% ===========================================
\usetikzlibrary{patterns,3d,backgrounds}
\usepgfplotslibrary{groupplots}

% Upgrade pgfplots compatibility
\pgfplotsset{compat=1.18}

% ===========================================
% OCEAN BLUE GRADIENT (p1-p9)
% ===========================================
\definecolor{p1}{HTML}{caf0f8}
\definecolor{p2}{HTML}{ade8f4}
\definecolor{p3}{HTML}{90e0ef}
\definecolor{p4}{HTML}{48cae4}
\definecolor{p5}{HTML}{00b4d8}
\definecolor{p6}{HTML}{0096c7}
\definecolor{p7}{HTML}{0077b6}
\definecolor{p8}{HTML}{023e8a}
\definecolor{p9}{HTML}{03045e}

% Dark background color
\definecolor{pag}{HTML}{293133}

% ===========================================
% MATERIAL DESIGN COLORS
% ===========================================
% Note: myred, myyellow already defined in main preamble
% These are additional/alternate versions
\definecolor{matred}{HTML}{f44336}
\definecolor{matpink}{HTML}{e81e63}
\definecolor{matpurple}{HTML}{9c27b0}
\definecolor{matdeeppurple}{HTML}{673ab7}
\definecolor{matindigo}{HTML}{3f51b5}
\definecolor{matblue}{HTML}{2196f3}
\definecolor{matlightblue}{HTML}{03a9f4}
\definecolor{matcyan}{HTML}{00bcd4}
\definecolor{matteal}{HTML}{009688}
\definecolor{matgreen}{HTML}{4caf50}
\definecolor{matlightgreen}{HTML}{8bc34a}
\definecolor{matlime}{HTML}{cddc39}
\definecolor{matyellow}{HTML}{ffeb3b}
\definecolor{matamber}{HTML}{ffc107}
\definecolor{matorange}{HTML}{ff9800}
\definecolor{matdeeporange}{HTML}{ff5722}
\definecolor{matbrown}{HTML}{795548}
\definecolor{matgray}{HTML}{9e9e9e}
\definecolor{matbluegray}{HTML}{607d8b}

% ===========================================
% COLORLET FOR TIKZ
% ===========================================
\colorlet{xcol}{blue!60!black}

% ===========================================
% DARK MODE TCOLORBOX
% ===========================================
\newtcolorbox{qq}[2][]{%
    boxrule=0.75pt,
    sharp corners,
    colframe=white,
    colback=pag,
    coltext=white,
    #1,
}

% Dark definition box
\newtcolorbox{darkdfn}[2][]{%
    enhanced,
    boxrule=1pt,
    sharp corners,
    colframe=p5,
    colback=pag,
    coltext=white,
    coltitle=white,
    fonttitle=\bfseries,
    title=#2,
    #1,
}

% Dark theorem box
\newtcolorbox{darkthm}[2][]{%
    enhanced,
    boxrule=1pt,
    sharp corners,
    colframe=matpurple,
    colback=pag,
    coltext=white,
    coltitle=white,
    fonttitle=\bfseries,
    title=#2,
    #1,
}

% ===========================================
% TIKZ 3D FIX
% ===========================================
% Small fix for canvas is xy plane at z
% https://tex.stackexchange.com/a/48776/121799
\makeatletter
\tikzoption{canvas is xy plane at z}[]{%
    \def\tikz@plane@origin{\pgfpointxyz{0}{0}{#1}}%
    \def\tikz@plane@x{\pgfpointxyz{1}{0}{#1}}%
    \def\tikz@plane@y{\pgfpointxyz{0}{1}{#1}}%
    \tikz@canvas@is@plane}
\makeatother

% ===========================================
% CONDITIONAL LINE TIKZSET
% ===========================================
\tikzset{
    conditional line/.style 2 args={
        /utils/exec={\pgfmathsetmacro{\xcoordA}{#1}},
        /utils/exec={\pgfmathsetmacro{\xcoordB}{#2}},
        insert path={
            \ifdim\xcoordA pt>\xcoordB pt
            (\xcoordA,0) -- (\xcoordB,0)
            \fi
        },
    },
}

% ===========================================
% PGFPLOTS DARK MODE DEFAULTS
% ===========================================
\pgfplotsset{
    dark axis/.style={
        axis line style={white},
        tick style={white},
        label style={white},
        grid style={pag!70},
        every axis label/.append style={white},
        every tick label/.append style={white},
    },
}

% ===========================================
% TIKZ EXTERNALIZATION (for fast compilation)
% ===========================================
% This creates .md5 cache files for figures
% Requires: pdflatex -shell-escape
%
% To enable, add AFTER loading this file:
%   \enabletikzexternalize
%
\newcommand{\enabletikzexternalize}{%
  \usetikzlibrary{external}%
  \tikzexternalize[prefix=figures/]%
}

% ===========================================
% CONVENIENT SHORTCUTS FOR DARK PLOTS
% ===========================================
\newcommand{\darkplot}[1]{%
    \begin{tikzpicture}
    \begin{axis}[
        dark axis,
        grid=both,
        major grid style={line width=.2pt,draw=pag!70},
        #1
    ]
}


% ===========================================
% QUICK GEOMETRY MACROS (tkz-euclide)
% ===========================================
% Draw a point: \point{name}{x}{y}
\newcommand{\gpoint}[3]{\tkzDefPoint(#2,#3){#1}}

% Draw line through points: \gline{A}{B}
\newcommand{\gline}[2]{\tkzDrawLine[color=p4](#1,#2)}

% Draw circle: \gcircle{center}{radius}
\newcommand{\gcircle}[2]{\tkzDrawCircle[color=p5](#1,#2)}

% Mark angle: \gangle{A}{B}{C}
\newcommand{\gangle}[3]{\tkzMarkAngle[color=p3,size=0.5](#1,#2,#3)}

% ===========================================
% USAGE EXAMPLE
% ===========================================
% In your document:
%
% \begin{tikzpicture}
%   \begin{axis}[dark axis, ...]
%     \addplot[p4, thick] {sin(deg(x))};
%   \end{axis}
% \end{tikzpicture}
%
% Or with qq box:
% \begin{qq}{My Title}
%   Content here in dark box
% \end{qq}


\course{Studies-in-Algebra-and-Group-Theory}

\begin{document}

\lecture{10}{Modules and Finitely Generated Abelian Groups}

We now take a detour from the purely linear-algebraic material to study modules---a generalization of vector spaces where the scalars come from a ring rather than a field. The payoff is immediate: the classification of finitely generated abelian groups reduces to ``linear algebra over \(\ZZ\).''

Recall from \textbf{Lecture 2} that every abelian group \((G, +)\) is a \(\ZZ\)-module: the action \(\ZZ \times G \to G\) sends \((n, g) \mapsto ng\) (where \(ng = g + g + \cdots + g\) for \(n > 0\), \(0 \cdot g = 0\), and \(ng = -(|n|g)\) for \(n < 0\)). This observation is the bridge between group theory and module theory.


\section{Modules over a Ring}

\subsection{Definitions and Examples}

\dfn{Module}{Let \(R\) be a ring (with unity \(1\)). A (left) \textbf{\(R\)-module} is an abelian group \((M, +)\) equipped with a scalar multiplication \(R \times M \to M\), \((r, m) \mapsto r \cdot m\), satisfying:
\begin{enumerate}[(i)]
    \item \(r(m_1 + m_2) = rm_1 + rm_2\) for all \(r \in R\), \(m_1, m_2 \in M\).
    \item \((r_1 + r_2)m = r_1 m + r_2 m\) for all \(r_1, r_2 \in R\), \(m \in M\).
    \item \((r_1 r_2)m = r_1(r_2 m)\) for all \(r_1, r_2 \in R\), \(m \in M\).
    \item \(1 \cdot m = m\) for all \(m \in M\).
\end{enumerate}
}

\ex{Vector spaces}{If \(R = k\) is a field, then a \(k\)-module is the same thing as a vector space over \(k\). So modules generalize vector spaces by allowing the scalars to come from a ring.
}

\ex{Abelian groups as \(\ZZ\)-modules}{Every abelian group \((G, +)\) is a \(\ZZ\)-module via \(n \cdot g = ng\). Conversely, every \(\ZZ\)-module is an abelian group (just forget the \(\ZZ\)-action, which is determined by the group structure anyway). The categories of abelian groups and \(\ZZ\)-modules are the same.
}

\ex{The ring as a module over itself}{\(R\) is an \(R\)-module via left multiplication: \(r \cdot s = rs\). This is a free module of rank 1.
}

\ex{\(R\)-modules from ring homomorphisms}{If \(\varphi : R \to S\) is a ring homomorphism, then any \(S\)-module \(M\) becomes an \(R\)-module via \(r \cdot m = \varphi(r) m\). This is called restriction of scalars.
}

\ex{Polynomial modules}{Let \(R = k[x]\) and \(V\) a vector space over \(k\). Any linear operator \(T : V \to V\) makes \(V\) into a \(k[x]\)-module via \(p(x) \cdot v = p(T)(v)\). Different operators give different module structures on the same underlying vector space.
}


\subsection{Submodules and Ideals}

\dfn{Submodule}{A \textbf{submodule} of an \(R\)-module \(M\) is a subset \(N \subseteq M\) that is itself an \(R\)-module under the inherited operations. Equivalently, \(N\) is a subgroup of \((M, +)\) that is closed under scalar multiplication: \(rn \in N\) for all \(r \in R\), \(n \in N\).
}

\dfn{Ideal}{A submodule of \(R\) (viewed as a module over itself) is called a \textbf{left ideal}. When \(R\) is commutative (as will be the case for us), left ideals are automatically \textbf{two-sided ideals}: subsets \(I \subseteq R\) satisfying:
\begin{enumerate}[(i)]
    \item \(I\) is an abelian subgroup of \((R, +)\).
    \item \(rI \subseteq I\) and \(Ir \subseteq I\) for all \(r \in R\).
\end{enumerate}
}

\ex{Ideals in \(\ZZ\) and \(k[x]\)}{The ideals of \(\ZZ\) are exactly the subgroups \(n\ZZ = \{nk : k \in \ZZ\}\) for \(n \geq 0\). Similarly, the ideals of \(k[x]\) are \(p(x) \cdot k[x] = \{p(x)q(x) : q(x) \in k[x]\}\). In both cases, every ideal is generated by a single element. Rings with this property are called \textbf{principal ideal domains} (PIDs). This is closely related to the Euclidean division algorithm: \(\operatorname{span}(p, q) = \operatorname{span}(\gcd(p, q))\).
}

\nt{The quotient of an \(R\)-module \(M\) by a submodule \(N\) is again an \(R\)-module, with \(r \cdot (m + N) = rm + N\). For example, \(\ZZ / n\ZZ = \ZZ/n\) as a \(\ZZ\)-module, and \(k[x] / (p(x)) \cdot k[x]\) is a \(k[x]\)-module. The quotient of \(R\) by an ideal is not just an \(R\)-module but a ring in its own right.}




\subsection{General Facts about Modules: Nothing Is True}

Over a field, every vector space has a basis, every linearly independent set extends to one, and every spanning set contains one. Over a general ring, all of these fail.

\ex{Bases need not exist}{The \(\ZZ\)-module \(M = \ZZ/n\) (for \(n \geq 2\)) has no basis. Any map \(\varphi : \ZZ \to M\) sending \(1 \mapsto m\) cannot be injective (since \(nm = 0\) in \(M\) but \(n \neq 0\) in \(\ZZ\)). So no single-element set is linearly independent, and no basis can exist.
}

\ex{Linearly independent sets may not be part of a basis}{Consider \(M = \ZZ\) as a \(\ZZ\)-module. The element \(2 \in \ZZ\) is linearly independent (if \(a \cdot 2 = 0\) then \(a = 0\)). But \(\{2\}\) cannot be extended to a basis of \(\ZZ\): the only bases of \(\ZZ\) are \(\{1\}\) and \(\{-1\}\), and no set containing \(2\) is a basis.
}

\ex{Spanning sets need not contain a basis}{The \(\ZZ\)-module \(M = \ZZ\) is spanned by \(\{4, 5\}\) (since \(\gcd(4, 5) = 1\), so \(1 = 5 \cdot 5 - 6 \cdot 4 \in \operatorname{span}(4, 5)\)). But \(\{4, 5\}\) is not linearly independent (\(5 \cdot 4 - 4 \cdot 5 = 0\)), and neither subset \(\{4\}\) or \(\{5\}\) spans \(\ZZ\). So the spanning set contains no basis.
}

\ex{Submodules of finitely generated modules need not be finitely generated}{Let \(R = k[x_1, x_2, \ldots]\) (polynomials in infinitely many variables) and \(M = R\) as an \(R\)-module. Then \(M\) is generated by \(\{1\}\). But the ideal \(M' = \{f \in R : f(0) = 0\}\) (polynomials with zero constant term) is a submodule not finitely generated: any finite subset only involves finitely many variables \(x_i\), so cannot span \(x_j\) for large \(j\).

By contrast, this pathology does not occur over Noetherian rings (including \(\ZZ\) and \(k[x_1, \ldots, x_n]\)).
}


\subsection{Module Homomorphisms}

\dfn{Module homomorphism}{Let \(M, N\) be \(R\)-modules. A map \(\varphi \in \Hom_R(M, N)\) is a module homomorphism if
\[\varphi(u + w) = \varphi(u) + \varphi(w) \quad \text{and} \quad \varphi(av) = a\varphi(v)\]
for all \(u, w \in M\), \(a \in R\).
}

\nt{When \(R\) is commutative, \(\Hom_R(M, N)\) is itself an \(R\)-module: \((\varphi + \psi)(v) = \varphi(v) + \psi(v)\) and \((a\varphi)(v) = a\varphi(v)\). (For noncommutative \(R\), this scalar multiplication need not yield an \(R\)-module homomorphism, so \(\Hom_R(M, N)\) is only an abelian group.) For free modules, things work as expected: \(\Hom_R(R^m, R^n) \cong R^{m \times n}\) (determined by images of basis vectors, just like linear maps between vector spaces). But for non-free modules, \(\Hom_R(M, N)\) can vanish even when both \(M\) and \(N\) are nonzero. For instance, \(\Hom_R(k, k[x]) = 0\) when \(R = k[x]\) and \(M = k = k[x]/(x)\): any \(\varphi\) must satisfy \(x \cdot \varphi(1) = \varphi(x \cdot 1) = \varphi(0) = 0\), so \(\varphi(1)\) is killed by \(x\) in \(k[x]\), forcing \(\varphi(1) = 0\).
}


\subsection{Free Modules and Direct Sums}

\dfn{Free module}{An \(R\)-module \(M\) is free if it has a basis: a set \(\{e_i\}_{i \in I}\) such that every element of \(M\) can be written uniquely as a finite \(R\)-linear combination \(\sum r_i e_i\). Equivalently, \(M \cong R^{(I)} = \bigoplus_{i \in I} R\). The rank of a free module is \(|I|\) (well-defined when \(R\) is commutative).%
}

\nt{Over a field, every module (= vector space) is free. Over \({\ZZ}\), free modules are the groups \({\ZZ}^n\). The module \({\ZZ}/n\) is \emph{not} free for \(n \geq 2\).}

\dfn{Direct sum of modules}{The direct sum \(M_1 \oplus M_2\) is the set of pairs \((m_1, m_2)\) with componentwise operations. More generally, \(\bigoplus_{i \in I} M_i\) consists of tuples \((m_i)_{i \in I}\) with only finitely many nonzero entries. This has the universal property: maps out of \(\bigoplus M_i\) correspond to families of maps out of each \(M_i\).%
}

\dfn{Exact sequences of modules}{A sequence of \(R\)-module homomorphisms
\[\cdots \to M_{i-1} \xrightarrow{f_{i-1}} M_i \xrightarrow{f_i} M_{i+1} \to \cdots\]
is exact at \(M_i\) if \(\Img(f_{i-1}) = \Ker(f_i)\). A short exact sequence \(0 \to A \xrightarrow{f} B \xrightarrow{g} C \to 0\) means \(f\) is injective, \(g\) is surjective, and \(\Img(f) = \Ker(g)\), i.e., \(C \cong B/f(A)\). A short exact sequence \emph{splits} if \(B \cong A \oplus C\). Over a PID, every short exact sequence of finitely generated modules with \(C\) free splits (analogous to Maschke's theorem).%
}


\section{Classification of Finitely Generated Abelian Groups}

The main theorem of this lecture is:

\thm{Classification of finitely generated abelian groups}{Every finitely generated abelian group \(G\) is isomorphic to a product of cyclic groups:
\[G \cong \ZZ/n_1 \times \ZZ/n_2 \times \cdots \times \ZZ/n_k \times \ZZ^\ell\]
where \(n_i \geq 2\) and \(\ell \geq 0\). Moreover, using the Chinese Remainder Theorem (\(\ZZ/mn \cong \ZZ/m \times \ZZ/n\) when \(\gcd(m, n) = 1\)), we can arrange each \(n_i\) to be a prime power. In this form, the decomposition is unique up to reordering.
}

The proof reduces to ``linear algebra over \(\ZZ\),'' following Artin \S14.4--14.7. The strategy has two ingredients: first, express any finitely generated \(\ZZ\)-module as \(\ZZ^n / \Img(T)\) for some \(T \in \Hom(\ZZ^m, \ZZ^n)\); second, diagonalize \(T\) by changes of basis over \(\ZZ\).


\subsection{Step 1: Presenting a Module as a Quotient}

\mprop{Presentation of finitely generated \(\ZZ\)-modules}{If \(M\) is a finitely generated \(\ZZ\)-module, then there exist \(m, n \geq 0\) and \(T \in \Hom(\ZZ^m, \ZZ^n)\) such that
\[M \cong \ZZ^n / \Img(T).\]
Equivalently, there is an exact sequence \(\ZZ^m \xrightarrow{T} \ZZ^n \to M \to 0\).
}

This says: quotient by a \(\ZZ\)-submodule of \(\ZZ^n\) is the same as quotient of an abelian group by a subgroup. We need a lemma:

\mlemma{Submodules of \(\ZZ^n\) are free}{Any submodule of \(\ZZ^n\) is finitely generated and free of rank \(\leq n\).}

\pf{Proof}{By induction on \(n\). For \(n = 1\): subgroups of \((\ZZ, +)\) are \(\{0\}\) or \(\ZZ a\) for some \(a \in \ZZ \setminus \{0\}\), both free of rank \(\leq 1\).

Assume the result for \(\ZZ^{n-1}\), and let \(N \subseteq \ZZ^n\) be a submodule. The projection \(\pi : \ZZ^n \to \ZZ^{n-1}\) sending \((a_1, \ldots, a_n) \mapsto (a_2, \ldots, a_n)\) restricts to a homomorphism \(\pi|_N : N \to \ZZ^{n-1}\).

The image \(\Img(\pi|_N)\) is a submodule of \(\ZZ^{n-1}\), hence free and finitely generated by induction. The kernel \(\Ker(\pi|_N) = N \cap (\ZZ \times 0 \times \cdots \times 0)\) is a subgroup of \(\ZZ\), hence free of rank 0 or 1.

If \(\Ker(\pi|_N)\) and \(\Img(\pi|_N)\) are both finitely generated (and free), then so is \(N\). Let \(e_1, \ldots, e_k\) be generators of \(\Ker(\pi|_N)\) (a basis) and \(g_1 = \pi(f_1), \ldots, g_m = \pi(f_m)\) generators of \(\Img(\pi|_N)\). For any \(x \in N\), write \(\pi(x) = \sum a_i g_i\), so \(\pi(x - \sum a_i f_i) = 0\), meaning \(x - \sum a_i f_i \in \Ker(\pi|_N) = \operatorname{span}(e_1, \ldots, e_k)\). Thus \(x \in \operatorname{span}(e_1, \ldots, e_k, f_1, \ldots, f_m)\), so \(\{e_i, f_j\}\) generates \(N\). That this set is a basis (i.e., free) follows from the directness of \(\Ker \cap \Img = 0\).
}

\pf{Proof of the proposition}{Let \(M\) be finitely generated with generators \(e_1, \ldots, e_n\). Define \(\varphi : \ZZ^n \to M\) by \((a_1, \ldots, a_n) \mapsto \sum a_i e_i\). This is surjective. The kernel \(N = \Ker(\varphi) \subseteq \ZZ^n\) is a submodule, hence finitely generated by the lemma. Let \(f_1, \ldots, f_m\) be generators of \(N\). Define \(T : \ZZ^m \to \ZZ^n\) by \(T(b_i) = f_i\). Then \(\Ker(\varphi) = \Img(T)\), giving the exact sequence
\[\ZZ^m \xrightarrow{T} \ZZ^n \xrightarrow{\varphi} M \to 0\]
and thus \(M \cong \ZZ^n / \Img(T)\).
}


\subsection{Step 2: Diagonalization over \(\ZZ\) (Smith Normal Form)}

The key difference from linear algebra over a field: we cannot divide freely in \(\ZZ\). The substitute is the Euclidean algorithm, which governs what ``row and column operations'' are available.

\dfn{Divisibility}{The \textbf{divisibility} of a nonzero element \(x = (a_1, \ldots, a_n) \in \ZZ^n\) is the largest \(d \in \ZZ_+\) such that \(x = dy\) for some \(y \in \ZZ^n\). Equivalently, \(d = \gcd(a_1, \ldots, a_n)\). An element is \textbf{primitive} if its divisibility is 1.
}

\mlemma{Primitive elements and bases}{An element of \(\ZZ^n\) can be part of a basis if and only if it is primitive. More generally, \(v \in \ZZ^n\) with divisibility \(d\) satisfies \(v = d \cdot w\) where \(w\) is a basis element.
}

\pf{Proof}{Elements of a basis \((e_1, \ldots, e_n)\) are primitive: linear independence prevents \(e_i = d(\sum a_j e_j)\) for \(d > 1\).

Conversely, let \(v = a_1 e_1 + \cdots + a_n e_n\) be primitive. Without loss of generality, \(a_1 \neq 0\) and \(|a_1| = \min\{|a_i| : a_i \neq 0\}\). For each \(k \geq 2\), write \(a_k = q_k a_1 + r_k\) by Euclidean division. Change basis to \((e_1' = e_1 + \sum_{k \geq 2} q_k e_k, \, e_2, \ldots, e_n)\). In this basis, \(v = a_1 e_1' + r_2 e_2 + \cdots + r_n e_n\) with \(0 \leq r_k < |a_1|\).

Repeating this process (which strictly decreases the minimum nonzero coefficient at each step) terminates in finitely many steps with \(v = d \cdot e_1''\) for some basis vector \(e_1''\). Since \(v\) is primitive, \(d = 1\).
}


\thm{Smith Normal Form}{For any \(T \in \Hom(\ZZ^m, \ZZ^n)\), there exist bases \((e_1, \ldots, e_m)\) of \(\ZZ^m\) and \((f_1, \ldots, f_n)\) of \(\ZZ^n\), and positive integers \(d_1, \ldots, d_r\) (where \(r \leq \min(m, n)\) is the rank of \(T\)), such that
\[T(e_i) = \begin{cases} d_i f_i & \text{if } 1 \leq i \leq r \\ 0 & \text{if } i > r. \end{cases}\]
In matrix form:
\[\mathcal{M}(T) = \begin{pmatrix} d_1 & & & 0 \\ & d_2 & & \\ & & \ddots & \\ & & & d_r & \\ 0 & & & & 0 \end{pmatrix}.\]
Moreover, the construction gives \(d_1 = \min\{\text{div}(T(x)) : x \notin \Ker(T)\} = \gcd\) of all entries of \(\mathcal{M}(T)\).
}

\pf{Proof}{By induction on \(m\).

If \(T = 0\), the statement is trivial for any \(m, n\).

\textbf{Case \(m = 1\):} \(T\) is determined by \(T(1) \in \ZZ^n\). Let \(d = \text{div}(T(1))\). By the lemma, \(T(1) = d f_1\) where \(f_1\) can be extended to a basis of \(\ZZ^n\). Take \(e_1 = 1\).

\textbf{Inductive step:} Assume \(T \neq 0\). Let \(d_1 = \min\{\text{div}(T(x)) : x \notin \Ker(T)\}\), and choose \(e_1\) with \(\text{div}(T(e_1)) = d_1\). Since \(T(e_1)\) has divisibility \(d_1\), write \(T(e_1) = d_1 f_1\) with \(f_1\) primitive, and extend \(f_1\) to a basis \((f_1, \ldots, f_n)\) of \(\ZZ^n\). Extend \(e_1\) to a basis \((e_1, \ldots, e_m)\) of \(\ZZ^m\) (using the lemma, since \(e_1\) is primitive: if \(e_1 = d \cdot e_1'\) then \(\text{div}(T(e_1')) = d_1/d\), contradicting minimality unless \(d = 1\)).

In the bases \((e_i)\) and \((f_j)\), the matrix of \(T\) has the form
\[\mathcal{M}(T) = \begin{pmatrix} d_1 & a_2 & \cdots & a_m \\ 0 & & & \\ \vdots & & \mathcal{M}(T') & \\ 0 & & & \end{pmatrix}\]
where \(T'\) is the restriction to \(\operatorname{span}(e_2, \ldots, e_m)\) composed with projection to \(\operatorname{span}(f_2, \ldots, f_n)\).

For \(j \geq 2\), write \(a_j = q_j d_1 + r_j\) with \(0 \leq r_j < d_1\). Change basis of \(\ZZ^m\) to \((e_1, e_2' = e_2 - q_2 e_1, \ldots, e_m' = e_m - q_m e_1)\). The matrix becomes
\[\begin{pmatrix} d_1 & r_2 & \cdots & r_m \\ 0 & & & \\ \vdots & & \mathcal{M}(T') & \\ 0 & & & \end{pmatrix}.\]
If any \(r_j \neq 0\), then \(\text{div}(T(e_j')) \leq r_j < d_1\), contradicting the minimality of \(d_1\). So \(r_j = 0\) for all \(j \geq 2\), and the matrix is
\[\begin{pmatrix} d_1 & 0 & \cdots & 0 \\ 0 & & & \\ \vdots & & \mathcal{M}(T') & \\ 0 & & & \end{pmatrix}.\]
By induction on \(m\), we can diagonalize \(T'\) by further basis changes in \(\operatorname{span}(e_2, \ldots, e_m)\) and \(\operatorname{span}(f_2, \ldots, f_n)\), giving the desired diagonal form.
}

\nt{The Smith Normal Form is the integer analogue of row reduction. Over a field, Gaussian elimination diagonalizes any matrix (in fact reduces it to have 1's and 0's on the diagonal). Over \(\ZZ\), the Euclidean algorithm plays the role of division, and the diagonal entries \(d_i\) need not be 1. The construction actually forces \(d_1 \mid d_2 \mid \cdots \mid d_r\) (the invariant factor divisibility chain), though we do not prove this here.}

\subsection{Step 3: Completing the Classification}

\pf{Proof of the classification theorem}{By Proposition 1 (presentation as quotient), any finitely generated abelian group \(M\) satisfies \(M \cong \ZZ^n / \Img(T)\) for some \(T \in \Hom(\ZZ^m, \ZZ^n)\). By the Smith Normal Form, after choosing appropriate bases, \(\Img(T)\) is spanned by \(d_1 f_1, \ldots, d_r f_r\) for some positive integers \(d_i > 0\) and \(r \leq n\). Therefore
\[M \cong \ZZ^n / \operatorname{span}(d_1 f_1, \ldots, d_r f_r) \cong \ZZ/d_1 \times \cdots \times \ZZ/d_r \times \ZZ^{n-r}.\]
The factors \(\ZZ/d_i\) with \(d_i = 1\) are trivial and can be dropped. Relabeling the remaining \(d_i \geq 2\) as \(n_1, \ldots, n_k\) and setting \(\ell = n - r\) gives the decomposition
\[M \cong \ZZ/n_1 \times \cdots \times \ZZ/n_k \times \ZZ^\ell.\]
By the Chinese Remainder Theorem, each \(\ZZ/n_i\) with \(n_i = p_1^{a_1} \cdots p_s^{a_s}\) decomposes as \(\ZZ/p_1^{a_1} \times \cdots \times \ZZ/p_s^{a_s}\). Rearranging gives the prime power form.
}

\ex{Finitely generated abelian groups of small order}{
\begin{itemize}
    \item Order 4: \(\ZZ/4\) or \(\ZZ/2 \times \ZZ/2\). These are not isomorphic (the first has an element of order 4, the second does not).
    \item Order 6: \(\ZZ/6 \cong \ZZ/2 \times \ZZ/3\) (only one, since \(\gcd(2,3) = 1\)).
    \item Order 8: \(\ZZ/8\), \(\ZZ/4 \times \ZZ/2\), or \(\ZZ/2 \times \ZZ/2 \times \ZZ/2\).
    \item Order 12: \(\ZZ/12 \cong \ZZ/4 \times \ZZ/3\), or \(\ZZ/2 \times \ZZ/6 \cong \ZZ/2 \times \ZZ/2 \times \ZZ/3\).
\end{itemize}
}

\nt{The number \(\ell\) is the \textbf{free rank} (or \textbf{Betti number}) of \(M\), and the integers \(n_1, \ldots, n_k\) (in prime power form) are the \textbf{elementary divisors}. These invariants completely determine \(M\) up to isomorphism. The free rank counts the number of ``copies of \(\ZZ\),'' while the elementary divisors encode the torsion (elements of finite order).}


\mer{Modules and finitely generated abelian groups exercises}{
\begin{enumerate}[(1)]
    \item Classify all abelian groups of order 36 up to isomorphism.
    
    \item Let \(T : \ZZ^2 \to \ZZ^2\) be the linear map with matrix \(\begin{pmatrix} 6 & 4 \\ 3 & 5 \end{pmatrix}\). Compute the Smith Normal Form of \(T\) and determine the group \(\ZZ^2 / \Img(T)\).
    
    \item Prove that a finitely generated abelian group is free (i.e., isomorphic to \(\ZZ^n\) for some \(n\)) if and only if it is torsion-free (i.e., the only element of finite order is 0).
    
    \item Let \(G\) be a finitely generated abelian group with \(|G| = p^n\) for a prime \(p\). Show that \(G \cong \ZZ/p^{a_1} \times \cdots \times \ZZ/p^{a_k}\) where \(a_1 + \cdots + a_k = n\) and \(a_i \geq 1\). How many such groups are there (as a function of the partition structure of \(n\))?
    
    \item Show that \({\ZZ}/m \otimes_{\ZZ} {\ZZ}/n \cong {\ZZ}/\gcd(m,n)\). (Hint: use the presentation \({\ZZ}/m = {\ZZ} / m{\ZZ}\) and the right-exactness of the tensor product.)
\end{enumerate}
}

\begin{solution}
\textbf{(1)} We have \(36 = 2^2 \cdot 3^2\). For the 2-part: \(\ZZ/4\) or \(\ZZ/2 \times \ZZ/2\). For the 3-part: \(\ZZ/9\) or \(\ZZ/3 \times \ZZ/3\). Taking all combinations:
\[\ZZ/4 \times \ZZ/9 \cong \ZZ/36, \quad \ZZ/4 \times \ZZ/3 \times \ZZ/3, \quad \ZZ/2 \times \ZZ/2 \times \ZZ/9, \quad \ZZ/2 \times \ZZ/2 \times \ZZ/3 \times \ZZ/3.\]
So there are exactly 4 abelian groups of order 36.

\textbf{(2)} Start with \(A = \begin{pmatrix} 6 & 4 \\ 3 & 5 \end{pmatrix}\). The gcd of all entries is 1, but let us apply row/column operations over \(\ZZ\).

\(R_1 \to R_1 - 2R_2\): \(\begin{pmatrix} 0 & -6 \\ 3 & 5 \end{pmatrix}\). Swap rows: \(\begin{pmatrix} 3 & 5 \\ 0 & -6 \end{pmatrix}\). \(C_2 \to C_2 - C_1\): \(\begin{pmatrix} 3 & 2 \\ 0 & -6 \end{pmatrix}\). \(C_1 \to C_1 - C_2\): \(\begin{pmatrix} 1 & 2 \\ 6 & -6 \end{pmatrix}\). \(C_2 \to C_2 - 2C_1\): \(\begin{pmatrix} 1 & 0 \\ 6 & -18 \end{pmatrix}\). \(R_2 \to R_2 - 6R_1\): \(\begin{pmatrix} 1 & 0 \\ 0 & -18 \end{pmatrix}\).

Smith Normal Form: \(\operatorname{diag}(1, 18)\). So \(\ZZ^2 / \Img(T) \cong \ZZ/1 \times \ZZ/18 \cong \ZZ/18 \cong \ZZ/2 \times \ZZ/9\).

Check: \(\det(A) = 30 - 12 = 18\), and \(|\ZZ^2 / \Img(T)| = |\det(T)| = 18\). Consistent.

\textbf{(3)} One direction: if \(G \cong \ZZ^n\), then \(G\) is torsion-free (only \(0 \in \ZZ^n\) has finite order).

Conversely, if \(G\) is finitely generated and torsion-free, write \(G \cong \ZZ/n_1 \times \cdots \times \ZZ/n_k \times \ZZ^\ell\) by the classification theorem. Each \(\ZZ/n_i\) contributes a nonzero element of finite order (namely \(\bar{1}\), which has order \(n_i\)). Since \(G\) is torsion-free, \(k = 0\), so \(G \cong \ZZ^\ell\).

\textbf{(4)} Since \(|G| = p^n\), the classification gives \(G \cong \ZZ/p^{a_1} \times \cdots \times \ZZ/p^{a_k}\) with \(a_i \geq 1\) and \(p^{a_1} \cdots p^{a_k} = p^n\), i.e., \(a_1 + \cdots + a_k = n\). The number of such groups equals the number of partitions of \(n\). For example, \(p^4\) gives partitions \((4), (3,1), (2,2), (2,1,1), (1,1,1,1)\), so 5 groups.

\textbf{(5)} From the short exact sequence \(0 \to \ZZ \xrightarrow{\times m} \ZZ \to \ZZ/m \to 0\), tensor with \(\ZZ/n\). Right-exactness gives:
\[\ZZ/n \xrightarrow{\times m} \ZZ/n \to \ZZ/m \otimes \ZZ/n \to 0.\]
The map \(\times m : \ZZ/n \to \ZZ/n\) has image \(m\ZZ/n\), which is the cyclic subgroup of order \(n/\gcd(m,n)\), i.e., \(m\ZZ/n \cong \ZZ/(n/\gcd(m,n))\). The quotient therefore has order \(\gcd(m,n)\):
\[\ZZ/m \otimes \ZZ/n \cong (\ZZ/n) / (m\ZZ/n) \cong \ZZ/\gcd(m,n).\]
\end{solution}


\end{document}